\clearpage
\SourceLeaf{486}%
\fancyhead[L]{\small BERICHTE ÜBER DIE GRADMESSUNG.}
\pdfbookmark[0]{Bericht über die Ergebnisse der trigonometrischen Messungen.}{bericht-resultate-1848}
\begin{center}\large [BERICHT ÜBER DIE ERGEBNISSE DER TRIGONOMETRISCHEN MESSUNGEN.]\end{center}
\par
— — —

Von den Rechnungen, mit denen aus dem Rohmaterial der Messungen [in den Journalen, die in der Reinschrift 35 Hefte füllen] die Endergebnisse [im Koordinatenverzeichnis] gewonnen wurden, habe ich dem Umfang nach nur einen sehr kleinen Teil in meine jetzigen Vorlagen aufnehmen können.

Wäre es in der Tat möglich, alle diese Rechnungen vollständig und in ganzer Ausführlichkeit wieder aufzustellen, so könnten sie leicht vier- oder sechsmal so viele Bände füllen [wie die Journale jetzt in gebundener Form ergeben]. Teils ist aber der größere Teil der Einzelheiten dieser Rechnungen überhaupt nicht aufbewahrt, teils befinden sich die noch immer sehr umfangreichen Bündel der aufbewahrten Papiere in einer Form, die eine vollständige und geordnete Wiederherstellung \textit{aller} Rechnungen fast darauf hinauslaufen ließe, meine gesamte Arbeit noch einmal zu wiederholen. Ich habe mich deshalb darauf beschränkt, den Teil der Zwischenrechnungen geordnet herauszuziehen, der als der aussagekräftigste und nützlichste gelten muss, nämlich die tabellarische Zusammenstellung aller an den verschiedenen Beobachtungsplätzen festgelegten und orientierten Richtungswinkel, wobei eine Parallele zum Göttinger Meridian den Null- oder Ausgangspunkt bildet. Diese tabellarischen Darstellungen sind unter der Bezeichnung \textit{Abrisse}\footnote{Die Gesamtzahl wird etwas über 450 betragen.} in sechs Heften zusammengefasst. Dabei hat mir Professor \textsc{Goldschmidt} in mehrfacher Hinsicht geholfen und zugleich die Reinschriften größtenteils selbst angefertigt. — — —

Zur Entschuldigung dafür, dass diese Arbeit so spät abgeschlossen wurde, muss ich noch bemerken, dass die Verspätung hauptsächlich dadurch entstand, dass zur Klärung der zahlreichen Zweifel, die sowohl bei den Abschriften als noch mehr bei der Anfertigung der Abrisse auftraten, nicht selten erst langwierige Nachforschungen nötig waren. — — —

Sobald es Zeit, Gesundheit und Kräfte erlauben, werden auf meine beiden ersten Abhandlungen über geodätische Probleme noch ein paar andere folgen, die näher auf die besonderen Gegenstände eingehen.

Göttingen, den 15. März 1848.
\begin{flushright}C. F. \textsc{Gauss}.\end{flushright}
